\subsection*{Висновки до третього розділу}
\addcontentsline{toc}{subsection}{Висновки до третього розділу}

На основі інтегративної теоретичної рамки (розділ~\ref{sec:1}) та емпіричних даних аналізу стану формування ІКК в українській дизайн-освіті (розділ~\ref{sec:2}) розроблено модель трансформації навчальних програм для формування ІКК майбутніх дизайнерів. Основні результати.

1. Розроблено \textbf{структурно-функціональну модель формування ІКК} майбутніх дизайнерів, що складається з 5 взаємопов'язаних блоків: мотиваційно-цільового (соціальне замовлення: 0\% покриття ШІ у 122 програмах, незадоволеність випускників до 33,3\%, рекомендація роботодавців щодо ШІ), теоретико-методологічного (TPACK + студійна педагогіка + компетентнісний підхід; 6 дидактичних принципів), змістового (4 компоненти ІКК, відображені на навчальні дисципліни), організаційно-технологічного (8 категорій ШІ-інструментів; модифіковані студійні методи) та діагностично-результативного (4 критерії; портфоліо-оцінювання; 4 рівні сформованості ІКК). Модель функціонує як цілісна система зі зворотним зв'язком.

2. Розроблено \textbf{рекомендації щодо трансформації навчальних дисциплін} для 6 типів курсів:
\begin{itemize}
    \item студійні курси (Живопис, Рисунок, Композиція): базовий рівень ШІ-інтеграції через метод <<ручне vs. ШІ>>; збереження фундаментальних навичок;
    \item курси цифрового дизайну (Adobe Suite, 3D, Figma, Моушн): інструментальний рівень через ШІ-аугментований робочий процес;
    \item теоретичні курси (Історія мистецтва, Дизайн-мислення): базовий рівень через критичний аналіз ШІ-контенту;
    \item професійні курси (UX/UI, Графічний дизайн, Ілюстрація): інструментальний/трансформаційний рівень через ШІ-інтегроване проєктне навчання;
    \item практики: трансформаційний рівень через вибір баз з ШІ-використанням та обов'язкову ШІ-документацію;
    \item загальноосвітні курси (Іноземна мова, Філософія): базовий рівень через ШІ-термінологію, ШІ-етику.
\end{itemize}

3. Розроблено \textbf{інтегровану модель 240-кредитного навчального плану} (4 роки, 8 семестрів) з наскрізною ШІ-інтеграцією. Ключові характеристики: 4 нові ШІ-дисципліни (18 кредитів, 7,5\%); 33 з 43 дисциплін (76,7\%) мають визначений рівень ШІ-інтеграції; <<ШІ-спіна>> від 1-го до 4-го курсу (<<Цифрова грамотність та основи AI>> $\to$ <<Генеративний дизайн та AI>> + <<Взаємодія людина-AI>> $\to$ <<Критичний дизайн та етика AI>>); прогресивне зростання складності від базового до трансформаційного рівня; збереження фундаментальних дисциплін (Рисунок, Живопис, Композиція).

4. Обґрунтовано \textbf{4 педагогічні умови} формування ІКК з емпіричною базою:
\begin{itemize}
    \item системна інтеграція ШІ наскрізно через навчальний план~--- відповідь на 0\% покриття ШІ та принцип TPACK;
    \item трансформація студійної педагогіки від <<виробника>> до <<виробника + куратора>>~--- відповідь на парадигмальний зсув, зафіксований у~\cite{Musiienko2026ScopingReview};
    \item узгодження з вимогами роботодавців~--- відповідь на пряму рекомендацію щодо ШІ та незадоволеність випускників;
    \item безперервний моніторинг ІКК через портфоліо-оцінювання~--- відповідь на недостатність рефлексивного компоненту.
\end{itemize}

5. Розроблено \textbf{стратегію поетапного впровадження} (4 етапи): Етап~1~--- усвідомлення та підготовка (семінари для викладачів, 1 нова дисципліна, пілотна ШІ-інтеграція у 3--4 дисципліни); Етап~2~--- інструментальна інтеграція (оновлення 8--10 силабусів, модифікація критики, портфоліо); Етап~3~--- поглиблена інтеграція (2 нові ШІ-дисципліни, ШІ-лабораторія, трансформаційний рівень); Етап~4~--- повна реалізація та оцінка (<<Критичний дизайн та етика AI>>, ШІ у дипломному проєкті, комплексна оцінка). Визначено ресурсне забезпечення: матеріально-технічне, кадрове та навчально-методичне.

6. Розроблено \textbf{дизайн експертної оцінки} моделі: опитувальник з 23 тверджень у 6 блоках (за 5-бальною шкалою), критерії відбору 5--7 експертів, методика обробки (середні значення, коефіцієнт конкордації Кендалла, тематичне кодування відкритих відповідей).

Запропонована модель є першою спробою системної трансформації навчальних програм спеціальності 022~Дизайн в Україні для формування ІКК з урахуванням ШІ-виміру. Вона ґрунтується на масштабній емпіричній базі (122 програми, 3\,043 освітні компоненти, 132 силабуси, опитування 32~випускників, 9~роботодавців, 61~студента, висновки 3~експертних груп НАЗЯВО, міжнародний бенчмаркінг 7~програм) та інтегративній теоретичній рамці~\cite{Musiienko2026ScopingReview, Musiienko2025DesignEducation}.
